\documentclass[11pt]{article}
\usepackage{amsmath,amssymb,amsfonts}
\usepackage{geometry}
\usepackage{hyperref}
\usepackage{graphicx}
\usepackage{physics}
\usepackage{bm}
\geometry{margin=1in}

\title{%
Prime Cantor Spacetime Algebra Stability Simulator\\[4pt]
\large A Defensive Publication on Prime–Indexed Cantor Embeddings,\\
Hausdorff Resonance at $2^{137}$, and Numerical Stability Metrics}
\author{Citizen Gardens \\ The Foundation of Multiplicity}
\date{April 27, 2026}

\begin{document}
\maketitle

\begin{abstract}
We document a numerical simulator and accompanying conceptual framework for Prime Cantor Spacetime Algebra (PCSA), in which spacetime embeddings are modeled by Cantor--type fractals indexed by primes.%
\footnote{For prior narrative and formal development of PCSA see \emph{Prime Cantor Spacetime Algebra (PCSA).}~\cite{pcsa_gist}. [web:21]} 
In this model, each prime $p$ defines a self--similar Cantor set with Hausdorff dimension
\[
  D_p = \frac{\ln\left(\frac{p+1}{2}\right)}{\ln p},
\]
and we focus on the dimension deficit $\delta_p = 1 - D_p$ as a proxy for spacetime embedding ``thickness.''%
\footnote{The same formula and its large--$p$ asymptotics are discussed in \cite{pcsa_gist}. [web:21]} 
For large $p$ near $2^{137}$, we find numerically that $\delta_p \approx 1/137$ up to double--precision roundoff, echoing the fine--structure constant $\alpha\approx 1/137$.%
\footnote{The resonance observation is articulated in \cite{pcsa_gist}. [web:21]} 

We build a Python simulator that (i) generates sampled Cantor--style fractal visualizations for selected Mersenne primes; (ii) computes Hausdorff dimensions and deficits for a family of primes; (iii) defines stability metrics such as a collapse index, resonance score, and mean gap to $1/137$; and (iv) produces static plots and an animated convergence diagram tracking the trade--off between resonance and collapse resistance when adding a fourth prime.%
\footnote{All numerical values and code snippets below are derived from the attached simulator implementation. [code\_file:43]} 
This report serves as a defensive publication: it fixes the constructions, numerical methodology, and quantitative claims in a citable form to establish prior art for prime--indexed Cantor spacetime embeddings and their use as stability diagnostics in multiplicity--based physics.
\end{abstract}
\newpage
\tableofcontents

\section{Executive Summary}

The work reported here lies at the intersection of number theory, fractal geometry, and speculative spacetime modeling. It extends a Prime Cantor Spacetime Algebra (PCSA) program in which the unit interval is partitioned by primes, and Cantor--style removal rules generate fractal sets that are interpreted as local spacetime fibers.%
\footnote{See the PCSA monograph for a broader conceptual framing of spacetime as a prime--indexed Cantor construction, including physical interpretations. [web:21]} 

\subsection{Core ideas}

\begin{itemize}
  \item \textbf{Prime--indexed Cantor sets.}
  For each prime $p$, we consider a self--similar Cantor set in $[0,1]$ whose Hausdorff dimension is
  \[
    D_p = \frac{\ln\left( \frac{p+1}{2} \right)}{\ln p}, \qquad 
    \delta_p := 1 - D_p
  \]
  with large--$p$ approximation $\delta_p \approx \ln 2 / \ln p$.%
  \footnote{This formula appears in PCSA as the dimension of a $p$--Cantor construction where $(p+1)/2$ subintervals are retained at each stage. [web:21]} 

  \item \textbf{Resonance at $p \approx 2^{137}$.}
  For $p$ close to $2^{137}$, the approximation $\delta_p \approx \ln 2 / \ln p$ yields $\delta_p \approx 1/137$, producing a numerical resonance with the fine--structure constant.%
  \footnote{The heuristic $1-D_p\approx\ln 2/\ln p$ and the choice $p\approx 2^{137}$ are highlighted in \cite{pcsa_gist}. [web:21]} 
  Our simulator confirms that for the exact anchor $p = 2^{137}$, $\delta_p$ equals $1/137$ to within $\sim 10^{-17}$, the scale of double--precision roundoff.%
  \footnote{The simulator computes $\delta_{2^{137}} \approx 0.007299270072992692$ with $|\delta_{2^{137}} - 1/137|\approx 8.67\times 10^{-18}$. [code\_file:43]} 

  \item \textbf{Visible versus extended prime sectors.}
  Motivated by PCSA's use of Mersenne primes as ``visible'' embedding channels, we study configurations based on exponents $31,61,89$ (3--prime sector) and $31,61,89,107$ (4--prime sector), with $p=2^k-1$ in each case.%
  \footnote{These exponents correspond to known Mersenne primes and are standard in the literature. [web:22][web:34]} 
  We treat the $2^{137}$ anchor as a resonance manifold rather than a visible prime.

  \item \textbf{Stability metrics.}
  We define:
  \begin{itemize}
    \item Embedding depth: the average Hausdorff dimension within a configuration.
    \item Collapse index: $1/\min_{i<j}|D_i - D_j|$, a proxy for singularity pressure from tightly clustered dimensions.
    \item Resonance score: an average of $1/(1+10^6|\delta_p - 1/137|)$ across primes in a configuration.
    \item Mean gap: the average $|\delta_p - 1/137|$ for primes in the configuration.
  \end{itemize}
  These metrics are used to quantify trade--offs between resonance and collapse resistance.%
  \footnote{All metric definitions are implemented explicitly in the simulator code. [code\_file:43]} 
\end{itemize}

\subsection{Numerical findings}

Direct evaluation shows the following representative values:%
\footnote{Numerical values here are produced by the Python code in Section~\ref{sec:code-snippet}. [code\_file:43]} 

\begin{itemize}
  \item For $M_{31},M_{61},M_{89}$ (3 primes): 
  embedding depth $\approx 0.9800$, collapse index $\approx 1.94\times 10^2$, resonance score $\approx 1.35\times 10^{-4}$, and mean gap $\approx 1.27\times 10^{-2}$.
  \item For $M_{31},M_{61},M_{89},M_{107}$ (4 primes):
  embedding depth $\approx 0.9827$, collapse index $\approx 5.29\times 10^2$, resonance score $\approx 2.23\times 10^{-4}$, and mean gap $\approx 1.00\times 10^{-2}$.
  \item For the $2^{137}$ anchor (and its nearest prime neighbor near $2^{137}$):
  embedding depth $\approx 0.9927007$, resonance score $\approx 1.0$, and mean gap $\approx 8.67\times 10^{-18}$.
\end{itemize}

Thus, adding a fourth prime improves resonance (smaller mean gap, larger resonance score) but increases the collapse index, indicating stronger singularity pressure from more tightly clustered dimensions.%
\footnote{The 4--prime collapse index is roughly $2.7\times$ the 3--prime value in the simulator output. [code\_file:43]} 
The $2^{137}$ anchor achieves near--perfect resonance but, as a single--prime manifold, is interpreted as a fixed embedding reference rather than a stable visible sector.

\section{Mathematical Overview}

\subsection{Prime Cantor Spacetime Algebra context}

PCSA models spacetime as a fractal construct built from prime--based Cantor sets, where primes label local ``fibers'' and their interactions.%
\footnote{A detailed exposition of PCSA and its physical motivation is presented in \cite{pcsa_gist}. [web:21]} 
The construction begins from the classical Cantor set (dimension $\ln 2/\ln 3$) and generalizes it to prime--indexed deletion rules; PCSA then interprets points in the resulting fractal as tuples encoding both combinatorial and physical data.%
\footnote{The PCSA paper relates these fractals to quantum vacuum structure and gauge transformations. [web:21]} 

\subsection{Hausdorff dimension of a $p$--Cantor set}

For fixed prime $p$, consider a self--similar Cantor construction that, at each step, divides each surviving interval into $p$ subintervals and retains $(p+1)/2$ of them.%
\footnote{This is the construction used in the PCSA monograph to derive $D_p$. [web:21]} 
The usual self--similarity argument for Hausdorff dimension gives
\[
  N r^{D_p} = 1,
\]
where $N = (p+1)/2$ is the number of retained subintervals and $r = 1/p$ is the contraction ratio per step. Solving for $D_p$ yields
\begin{align}
  D_p 
    &= \frac{\ln N}{\ln(1/r)} 
     = \frac{\ln((p+1)/2)}{\ln p}. \label{eq:Dp}
\end{align}
For large primes $p$, one obtains the asymptotic expansion
\begin{align}
  D_p 
    &= \frac{\ln p + \ln\left(\frac{1 + 1/p}{2}\right)}{\ln p} \\
    &\approx 1 - \frac{\ln 2}{\ln p} + O\!\left(\frac{1}{p \ln p}\right),
\end{align}
so the deficit
\begin{equation}
  \delta_p := 1 - D_p \approx \frac{\ln 2}{\ln p}.
\end{equation}

The formalism of $p$--Cantor sets and their Hausdorff measures is part of a broader literature on generalized Cantor sets and $h$--Hausdorff measures.%
\footnote{See, for instance, work on $p$--Cantor sets and their Hausdorff measure such as \cite{hausdorff_pcantor}. [web:31][web:61]} 
Prime--restricted Cantor sets based on continued fractions also exhibit nontrivial dimension phenomena.%
\footnote{For a modern treatment of prime continued--fraction Cantor sets and their exact Hausdorff and packing dimensions, see Das and Simmons~\cite{das_simmons_prime_cantor}. [web:39]} 

\subsection{Resonance at $p \approx 2^{137}$}

Using the asymptotic $\delta_p \approx \ln 2 / \ln p$, choosing $p \approx 2^{137}$ gives
\[
  \delta_p \approx \frac{\ln 2}{\ln 2^{137}} 
  = \frac{\ln 2}{137\ln 2} 
  = \frac{1}{137},
\]
which aligns numerically with the fine--structure constant $\alpha\approx 1/137$.%
\footnote{This numerical coincidence is one of the motivating observations in the PCSA framework. [web:21]} 
Our simulator evaluates $\delta_p$ for the exact anchor $p=2^{137}$ and finds
\begin{align}
  \delta_{2^{137}} 
    &= 1 - \frac{\ln\bigl((2^{137}+1)/2\bigr)}{\ln(2^{137})} \\
    &\approx 0.007299270072992692,
\end{align}
with 
\[
  \Bigl|\delta_{2^{137}} - \frac{1}{137}\Bigr| \approx 8.67\times 10^{-18},
\]
which is at the level of float64 machine epsilon for numbers of this magnitude.%
\footnote{These values are directly output by the simulator; see Section~\ref{sec:code-snippet}. [code\_file:43]} 
In the PCSA interpretation, $2^{137}$ is therefore a \emph{resonance anchor}: a prime scale at which the combinatorial thickness of the prime Cantor set matches a fundamental physical constant to numerical precision.

\subsection{Configurations based on Mersenne primes}

To model ``visible'' primes, we restrict attention to Mersenne primes $M_k = 2^k-1$ with known exponents such as $k=31,61,89,107$.%
\footnote{These exponents are well--known entries in the list of Mersenne primes. [web:22][web:34]} 
For each $M_k$ we compute $D_{M_k}$ via~\eqref{eq:Dp} and $\delta_{M_k}=1-D_{M_k}$, and we study configurations:
\[
  \mathcal{C}_3 = \{M_{31},M_{61},M_{89}\},\quad
  \mathcal{C}_4 = \{M_{31},M_{61},M_{89},M_{107}\},
\]
with a separate anchor configuration containing only $2^{137}$.%

\section{Stability Metrics}

We now formalize the stability metrics implemented in the simulator.

\subsection{Dimension and deficit}

For each prime $p$ we define
\begin{equation}
  D_p = \frac{\ln((p+1)/2)}{\ln p}, \qquad
  \delta_p = 1 - D_p.
\end{equation}
Given a target value $\delta_\ast = 1/137$, we consider the gap
\begin{equation}
  g(p) := \abs{\delta_p - \delta_\ast}.
\end{equation}

\subsection{Embedding depth}

For a configuration $\mathcal{C} = \{p_1,\dots,p_n\}$, we define the embedding depth as the mean Hausdorff dimension
\begin{equation}
  \mathrm{depth}(\mathcal{C}) 
    := \frac{1}{n}\sum_{i=1}^n D_{p_i}.
\end{equation}
In the simulator output, $\mathcal{C}_3$ has depth $\approx 0.9800$ while $\mathcal{C}_4$ has depth $\approx 0.9827$, and the $2^{137}$ anchor alone has depth $\approx 0.9927$.%
\footnote{These averaged values come from the computed $D_p$ values for each configuration. [code\_file:43]} 

\subsection{Collapse index}

To estimate a kind of singularity pressure associated with tightly clustered dimensions, we define
\begin{equation}
  \mathrm{collapse}(\mathcal{C}) 
    := \frac{1}{\min_{i<j}|D_{p_i} - D_{p_j}|},
\end{equation}
with an $\varepsilon$--floor to avoid division by zero numerically. Configurations with more tightly packed dimensions have smaller minimum separations and thus larger collapse indices.

In the simulator:
\begin{itemize}
  \item $\mathrm{collapse}(\mathcal{C}_3) \approx 1.94\times 10^2$.
  \item $\mathrm{collapse}(\mathcal{C}_4) \approx 5.29\times 10^2$.
\end{itemize}
Thus, adding the fourth prime increases the collapse index by a factor of about $2.7$, reflecting increased singularity pressure under this measure.%
\footnote{Values from the summary CSV generated by the simulator. [code\_file:43]} 

\subsection{Resonance score}

To capture how well a configuration matches the target deficit $\delta_\ast$, we define
\begin{equation}
  \mathrm{resonance}(\mathcal{C}) 
    := \frac{1}{n}\sum_{i=1}^n \frac{1}{1 + 10^6\, g(p_i)}.
\end{equation}
This maps exact matches ($g(p_i)=0$) to score $1$ and strongly penalizes large gaps. In the simulator:
\begin{itemize}
  \item $\mathrm{resonance}(\mathcal{C}_3) \approx 1.35\times 10^{-4}$.
  \item $\mathrm{resonance}(\mathcal{C}_4) \approx 2.23\times 10^{-4}$.
  \item $\mathrm{resonance}(\{2^{137}\}) \approx 1.0$ (within machine precision).
\end{itemize}
Hence, the 4--prime configuration is more resonant than the 3--prime sector, while the $2^{137}$ anchor is essentially perfectly resonant.%
\footnote{Again, these values are taken from the simulator’s summary. [code\_file:43]} 

\subsection{Mean gap}

Finally, we track the mean absolute gap
\begin{equation}
  \overline{g}(\mathcal{C}) := \frac{1}{n}\sum_{i=1}^n g(p_i).
\end{equation}
Typical values are:
\begin{itemize}
  \item $\overline{g}(\mathcal{C}_3) \approx 1.27\times 10^{-2}$.
  \item $\overline{g}(\mathcal{C}_4) \approx 1.00\times 10^{-2}$.
  \item $\overline{g}(\{2^{137}\}) \approx 8.67\times 10^{-18}$.
\end{itemize}
These confirm that adding the fourth prime brings the visible sector closer to the $1/137$ resonance while the anchor lies essentially exactly on it.%
\footnote{Mean gap values are directly read from the simulator’s per--configuration statistics. [code\_file:43]} 

\section{Simulator Implementation}
\label{sec:code-snippet}

This section documents the core Python implementation of the PCSA stability simulator. The design goals are:

\begin{itemize}
  \item Compute $D_p$, $\delta_p$, and $g(p)$ for a selected family of primes.
  \item Construct basic stability metrics for configurations of primes.
  \item Generate sampled Cantor--style fractal visualizations.
  \item Produce convergence plots and an animated sweep for the fourth prime.
\end{itemize}

\subsection{Core numerical kernel}

The following code snippet, extracted and slightly condensed from the full implementation, shows the central functions for computing $D_p$, $\delta_p$, and configuration metrics.%
\footnote{The full simulator script is provided as \texttt{pcsa\_simulator.py} in the artifacts. [code\_file:43]} 

\begin{verbatim}
import math
import numpy as np

def D_p(p: int) -> float:
    """Hausdorff dimension D_p = ln((p+1)/2) / ln p."""
    return math.log((p + 1) / 2.0) / math.log(p)

def delta_p(p: int) -> float:
    """Dimension deficit delta_p = 1 - D_p."""
    return 1.0 - D_p(p)

def gap_to_target(p: int, target: float = 1/137) -> float:
    """Absolute gap |delta_p - target|."""
    return abs(delta_p(p) - target)

def embedding_depth(primes) -> float:
    """Mean Hausdorff dimension over a configuration."""
    Ds = [D_p(p) for p in primes]
    return float(np.mean(Ds))

def collapse_index(primes, eps: float = 1e-18) -> float:
    """1 / min_{i<j} |D_i - D_j|, with epsilon floor."""
    Ds = [D_p(p) for p in primes]
    gaps = []
    for i in range(len(Ds)):
        for j in range(i + 1, len(Ds)):
            gaps.append(abs(Ds[i] - Ds[j]))
    return 1.0 / max(min(gaps), eps) if gaps else float('nan')

def resonance_score(primes, target: float = 1/137, scale: float = 1e6) -> float:
    """Average of 1 / (1 + scale * |delta_p - target|)."""
    vals = [1.0 / (1.0 + scale * gap_to_target(p, target)) for p in primes]
    return float(np.mean(vals))

def mean_gap(primes, target: float = 1/137) -> float:
    """Mean absolute gap to the target deficit."""
    return float(np.mean([gap_to_target(p, target) for p in primes]))
\end{verbatim}

This kernel abstracts away from any particular choice of primes; a separate layer specifies Mersenne exponents and the $2^{137}$ anchor.

\subsection{Prime choices and configurations}

In the current instantiation, we choose:
\begin{itemize}
  \item Visible Mersenne exponents: $k \in \{31,61,89\}$.
  \item Fourth prime exponent: $k=107$.
  \item Anchor exponent: $k=137$ with $p = 2^{137}$.
\end{itemize}
The corresponding primes are
\[
  M_k = 2^k - 1,\quad p_\text{anchor} = 2^{137}.
\]

\subsection{Fractal visualization}

Exact interval enumeration for large primes is computationally infeasible, so the simulator uses a sampled visualization strategy: at level $\ell$, it draws a finite number of uniformly spaced subintervals of nominal length $\sim p^{-\ell}$ to suggest the Cantor structure.%
\footnote{This design choice is made purely for numerical tractability; the underlying self--similar structure is defined analytically by the Hausdorff dimension. [code\_file:43]} 
An illustrative function is:

\begin{verbatim}
from matplotlib.patches import Rectangle

def sample_level_segments(p: int, level: int, max_segments: int = 180):
    """Sampled subintervals at a given iteration level."""
    n = min(max_segments, max(1, 2**level))
    seg_len = max(1.0 / (p ** level), 0.0015)
    starts = np.linspace(0.0, max(0.0, 1.0 - seg_len), n)
    return [(float(s), float(min(1.0, s + seg_len))) for s in starts]
\end{verbatim}

These segments are rendered as horizontal bars stacked across iterations and primes, producing a multi--row figure where each row corresponds to a chosen Mersenne prime and each column to an iteration.%
\footnote{The simulator saves this figure as \texttt{pcsa\_fractal\_iterations.png}. [code\_file:43]} 

\subsection{Convergence sweep and animation}

To illustrate how $\delta_p$ converges toward $1/137$ as the exponent $k$ grows, the simulator computes
\[
  \delta_{2^k} = 1 - \frac{\ln\left((2^k+1)/2\right)}{\ln(2^k)}
\]
for $k$ from $2$ to $200$, and plots both $\delta_{2^k}$ and $|\delta_{2^k} - 1/137|$ (the latter on a log scale).%
\footnote{This sweep is saved as \texttt{pcsa\_delta\_sweep.png}. [code\_file:43]} 
An animated sweep then varies the candidate fourth--prime exponent $k$ across a range, plotting the instantaneous $\delta_{2^k}$ and updating bar plots of embedding depth and resonance score for the 3--prime versus 4--prime configurations.%
\footnote{The resulting GIF is stored as \texttt{pcsa\_convergence.gif}. [code\_file:43]} 

A minimal skeleton of the animation loop is:

\begin{verbatim}
from matplotlib.animation import FuncAnimation, PillowWriter

ks = np.arange(2, 201)
deltas = np.array([delta_p(2**int(k)) for k in ks], dtype=float)

def update(k):
    # Clear and redraw both subplots for current k
    # (omitted here: plotting calls and labels)
    return []

fig, axs = plt.subplots(1, 2, figsize=(12, 5))
frames = list(range(2, 138, 4))
ani = FuncAnimation(fig, update, frames=frames, interval=150, blit=False)
ani.save("pcsa_convergence.gif", writer=PillowWriter(fps=6))
\end{verbatim}

\section{Prior Art and Related Work}

\subsection{Prime Cantor Spacetime Algebra and Breathing Prime Cosmology}

The present simulator is best understood as a numerical module within a larger theoretical program that includes:
\begin{itemize}
  \item A Prime Cantor Spacetime Algebra, modeling spacetime as a prime--indexed Cantor set with an associated noncommutative algebra of ``points.''%
  \footnote{See the PCSA monograph for definitions of fractal points, non--associative addition, and collapse operations. [web:21]} 
  \item Breathing Prime Cosmology, which treats prime distributions as inducing curvature in a number--theoretic ``spacetime,'' linking global zeta geometry to local prime fibers.%
  \footnote{A unified exposition of Breathing Prime Cosmology and PCSA is given in the ``Breathing Prime Universe'' monograph. [web:24]} 
\end{itemize}
These works already articulate the idea of using prime--based fractals as models of quantum vacuum structure and the possibility of matching fractal dimension deficits to physical constants such as $\alpha$.%
\footnote{The link between $1-D_p$ and the fine--structure constant is emphasized as a heuristic resonance in \cite{pcsa_gist}. [web:21]} 

\subsection{Fractal Cantor sets and Hausdorff dimension}

The classical Cantor set and its Hausdorff dimension $\ln 2/\ln 3$ are well established.%
\footnote{Standard expositions of Cantor sets and their dimension can be found in introductory fractal geometry texts and expository articles. [web:51]} 
There is extensive literature on generalized and $p$--Cantor sets and the behavior of their Hausdorff and $h$--Hausdorff measures.%
\footnote{For example, see work on $p$--Cantor sets and their $h$--Hausdorff measures. [web:31][web:58][web:61]} 
More recently, prime--restricted Cantor sets defined via continued fractions have been studied, with exact Hausdorff and packing dimensions determined under mild number--theoretic assumptions.%
\footnote{See Das and Simmons, ``Exact dimensions of the prime continued fraction Cantor set,'' arXiv:2305.11829. [web:39][web:62]} 

The PCSA framework is distinct from these: it uses a simple closed--form dimension formula for a particular prime--indexed Cantor construction and interprets it physically in terms of spacetime embedding and resonance.

\subsection{Spacetime algebra and geometric approaches}

Spacetime algebra in the sense of Clifford/Geometric Algebra has been used extensively in relativistic physics and quantum information to provide coordinate--free formulations of field equations and quantum gates.%
\footnote{See, e.g., standard references on spacetime algebra and its use in quantum computing and multiparticle spacetime algebra. [web:56][web:57]} 
PCSA differs by using prime--indexed fractal geometry as its organizing principle rather than a Clifford algebra of Minkowski space, though conceptually both approaches aim at unification of algebraic and geometric structure.

\subsection{Multiplicity and multiset perspectives}

The broader idea of treating sets as multiplicities---tallies of repeated elements---has its roots in multiset theory and Cantor's own notions of multiplicity.%
\footnote{For historical context on multiplicities and multisets, see, for example, early work on multiset theory and Cantor's usage. [web:60]} 
Multiplicity Theory, in the present context, is extended into a prime--indexed, recursively stable framework where multiplicity spaces become prime--labeled fractal structures rather than mere collections.

\section{Claim Scope and Defensive Publication Notes}

The constructions documented here support the following non--exclusive claims for prior art:

\begin{itemize}
  \item Use of prime--indexed Cantor sets with Hausdorff dimension $D_p = \ln((p+1)/2)/\ln p$ as models of spacetime embedding layers and their thickness via $\delta_p = 1-D_p$.%
  \footnote{The specific formula, interpretation, and use of $\delta_p$ as an embedding deficit come from the PCSA framework. [web:21]} 
  \item Identification of $p\approx 2^{137}$ and in particular $p = 2^{137}$ as a resonance anchor where $\delta_p \approx 1/137$, with explicit numerical confirmation of this equality to float64 precision.%
  \footnote{The simulator furnishes concrete numerical evidence of this resonance. [code\_file:43]} 
  \item Definition of stability metrics (embedding depth, collapse index, resonance score, mean gap) applied to configurations of primes (e.g., Mersenne prime sectors) to quantify trade--offs between resonance with $\alpha$ and singularity pressure.
  \item Construction of a numerical simulator that:
    \begin{itemize}
      \item Implements the above metrics.
      \item Generates sampled Cantor--style visualizations for large primes.
      \item Produces static and animated plots capturing convergence toward $1/137$.
    \end{itemize}
\end{itemize}

As such, this document, together with the accompanying code and generated artifacts, acts as a defensive publication for any application of prime--indexed Cantor fractals with the specific $D_p$ above as spacetime embedding models whose thickness resonates with the fine--structure constant via $p\approx 2^{137}$, and for numerical stability analyses based on these metrics.

\section*{Acknowledgments and Future Work}

This work builds on the PCSA and Breathing Prime Cosmology monographs that first articulated the prime--Cantor spacetime viewpoint and its resonance with $1/137$.%
\footnote{See \cite{pcsa_gist} and \cite{breathing_prime}. [web:21][web:24]} 
Future directions include:
\begin{itemize}
  \item Refining the collapse index into an energy--weighted or measure--theoretic functional.
  \item Extending the simulator to configurations that use prime sequences beyond Mersenne primes.
  \item Integrating the PCSA simulator into a broader multiplicity--theoretic toolkit for modeling nonlinear, emergent structures across physics and cognition.
\end{itemize}

\appendix
\section*{Mathematical Appendix: Dimension, Resonance, and Operator Bounds}

\section{Proof of the Prime–Indexed Cantor Dimension Formula}

We recall the specific prime–indexed Cantor construction used throughout the article, following the Prime Cantor Spacetime Algebra (PCSA) exposition.%
\footnote{The same construction and resulting dimension formula are stated in \cite{pcsa_gist}. [web:21]} 

\subsection{Self–similar construction}

Fix a prime \(p\in \mathbb{N}\). Consider the following self–similar subset \(C_p \subset [0,1]\):
\begin{enumerate}
  \item Start from \(C_p^{(0)} = [0,1]\).
  \item At each step, replace every interval in \(C_p^{(n)}\) by \(p\) equal subintervals of length \(r = 1/p\).
  \item Retain exactly \(N = \frac{p+1}{2}\) of these subintervals according to a parity or residue rule; delete the rest.%
  \footnote{PCSA uses a prime–parity rule for retention, but for the dimension calculation only the count \(N\) and ratio \(r\) matter. [web:21]} 
\end{enumerate}
Define \(C_p = \bigcap_{n\ge 0} C_p^{(n)}\). Since the construction is self–similar with \(N\) subintervals of identical contraction ratio \(r\) at each step and satisfies an open set condition, the Hausdorff dimension of \(C_p\) is determined by the Moran equation.%
\footnote{For self–similar sets satisfying the open set condition, the Hausdorff dimension is the unique solution \(d\) to \(\sum_i r_i^{d}=1\); see standard references on Hausdorff dimension and Moran's theorem. [web:33][web:75]} 

\subsection{Proof of the dimension formula}

\begin{proposition}
Let \(C_p\) be the self–similar set defined above. Then its Hausdorff dimension \(D_p\) is given by
\[
  D_p = \frac{\ln\left(\frac{p+1}{2}\right)}{\ln p}.
\]
\end{proposition}

\begin{proof}
At each stage, the set \(C_p\) is the union of \(N = \frac{p+1}{2}\) pairwise disjoint scaled copies of itself, each obtained by an affine contraction of ratio \(r = 1/p\). By self–similarity under the open set condition, the Hausdorff dimension \(D_p\) must satisfy the Moran (similarity) equation
\[
  \sum_{i=1}^{N} r^{D_p} = 1.
\]
Since all contraction ratios are equal to \(r\), this reduces to
\[
  N r^{D_p} = 1
  \quad\Longleftrightarrow\quad
  \left(\frac{p+1}{2}\right)\left(\frac{1}{p}\right)^{D_p} = 1.
\]
Taking natural logarithms of both sides,
\[
  \ln\left(\frac{p+1}{2}\right) + D_p \ln\left(\frac{1}{p}\right) = 0.
\]
Since \(\ln(1/p) = -\ln p\), this becomes
\[
  \ln\left(\frac{p+1}{2}\right) - D_p \ln p = 0
  \quad\Longleftrightarrow\quad
  D_p = \frac{\ln\left(\frac{p+1}{2}\right)}{\ln p},
\]
which is the desired formula.
\end{proof}

\subsection{Asymptotics of the dimension deficit}

Define the deficit
\[
  \delta_p := 1 - D_p.
\]
We obtain the large–\(p\) asymptotics that drive the resonance analysis.

\begin{lemma}[Asymptotic expansion of \(\delta_p\)]
For primes \(p\) tending to infinity,
\[
  \delta_p = 1 - \frac{\ln\left(\frac{p+1}{2}\right)}{\ln p}
  = \frac{\ln 2}{\ln p} + O\!\left(\frac{1}{p \ln p}\right).
\]
\end{lemma}

\begin{proof}
Write
\[
  \ln\left(\frac{p+1}{2}\right)
   = \ln p + \ln\left(\frac{1 + 1/p}{2}\right)
   = \ln p + \ln(1 + 1/p) - \ln 2.
\]
Using the Taylor expansion \(\ln(1 + 1/p) = \frac{1}{p} + O(1/p^2)\), we have
\[
  \ln\left(\frac{p+1}{2}\right)
  = \ln p - \ln 2 + \frac{1}{p} + O\!\left(\frac{1}{p^2}\right).
\]
Therefore
\[
  D_p = \frac{\ln\left(\frac{p+1}{2}\right)}{\ln p}
      = \frac{\ln p - \ln 2}{\ln p} + \frac{1}{p\ln p} + O\!\left(\frac{1}{p^2\ln p}\right)
      = 1 - \frac{\ln 2}{\ln p} + O\!\left(\frac{1}{p\ln p}\right).
\]
Subtracting from \(1\) yields
\[
  \delta_p = 1 - D_p = \frac{\ln 2}{\ln p} + O\!\left(\frac{1}{p\ln p}\right),
\]
as claimed.
\end{proof}

\section{Resonance at the \(2^{137}\) Anchor}

\subsection{Exact evaluation at \(p = 2^{137}\)}

\begin{proposition}[Resonance identity at \(2^{137}\)]
Let \(p = 2^{137}\). Then the deficit
\[
  \delta_p = 1 - \frac{\ln\left(\frac{p+1}{2}\right)}{\ln p}
\]
satisfies
\[
  \delta_p = \frac{1}{137} + \varepsilon_p,
\]
where \(|\varepsilon_p|\) is smaller than or comparable to the rounding scale of double–precision floating point arithmetic used in the simulator (empirically \(|\varepsilon_p|\approx 8.67\times 10^{-18}\)).%
\footnote{The numerical value of \(\varepsilon_p\) is obtained from the simulator; see the main text and code outputs. [code\_file:43]} 
\end{proposition}

\begin{proof}[Proof sketch]
For \(p = 2^{137}\),
\[
  D_p = \frac{\ln\left(\frac{2^{137}+1}{2}\right)}{\ln(2^{137})}
      = \frac{\ln\bigl(2^{136} + \tfrac{1}{2}\bigr)}{137\ln 2}.
\]
Write
\[
  \ln\bigl(2^{136} + \tfrac{1}{2}\bigr)
    = \ln\left(2^{136}\left(1 + \frac{1}{2\cdot 2^{136}}\right)\right)
    = 136\ln 2 + \ln\left(1 + \frac{1}{2^{137}}\right).
\]
Thus
\[
  D_p = \frac{136\ln 2 + \ln\bigl(1 + 2^{-137}\bigr)}{137\ln 2}
      = \frac{136}{137} + \frac{1}{137\ln 2}\,\ln\bigl(1 + 2^{-137}\bigr).
\]
Consequently
\[
  \delta_p 
    = 1 - D_p
    = 1 - \frac{136}{137} - \frac{1}{137\ln 2}\,\ln\bigl(1 + 2^{-137}\bigr)
    = \frac{1}{137} - \frac{1}{137\ln 2}\,\ln\bigl(1 + 2^{-137}\bigr).
\]
Using \(\ln(1+x) = x + O(x^2)\) with \(x = 2^{-137}\), we have
\[
  \ln\bigl(1 + 2^{-137}\bigr) = 2^{-137} + O(2^{-274}),
\]
so that
\[
  \delta_p
    = \frac{1}{137} - \frac{2^{-137}}{137\ln 2} + O\!\left(\frac{2^{-274}}{\ln 2}\right)
    = \frac{1}{137} + \varepsilon_p,
\]
with
\[
  \varepsilon_p := - \frac{2^{-137}}{137\ln 2} + O(2^{-274}).
\]
The leading term has magnitude
\[
  \left|\frac{2^{-137}}{137\ln 2}\right| \ll 2^{-137},
\]
which is extraordinarily small (\(\sim 10^{-41}\)). In numerical practice, the dominant discrepancy observed in the simulator (\(\approx 8.67\times 10^{-18}\)) arises from floating–point rounding in evaluating the logarithms and division, which dwarfs the analytic error term.%
\footnote{The simulator reports \(|\delta_{2^{137}} - 1/137| \approx 8.67\times 10^{-18}\), interpretable as numerical roundoff; see the CSV outputs. [code\_file:43]} 
\end{proof}

\subsection{Convergence along powers of two}

For general exponents \(k\), define
\[
  \delta_{2^k} := 1 - \frac{\ln\left(\frac{2^{k}+1}{2}\right)}{\ln(2^k)}.
\]
Using the same technique as above, we obtain
\[
  \delta_{2^k}
   = \frac{1}{k} - \frac{2^{-k}}{k\ln 2} + O(2^{-2k}),
\]
so that
\[
  \delta_{2^k} = \frac{1}{k} + O(2^{-k}).
\]
In particular,
\[
  \biggl|\delta_{2^k} - \frac{1}{k}\biggr| \le \frac{C}{k} 2^{-k}
\]
for some absolute constant \(C>0\), giving a rapidly decaying bound on the deviation from the simple \(1/k\) heuristic.

\section{Operator–Norm Style Bounds for Stability Metrics}

We now formalize operator–norm style bounds for the stability metrics treated numerically in the simulator. These bounds are simple but explicit, and they provide analytical control over the functionals used in the PCSA stability analysis.

\subsection{Configuration as a vector of dimensions}

Let \(\mathcal{C} = \{p_1,\dots,p_n\}\) be a finite configuration of primes. Define
\[
  \mathbf{D}(\mathcal{C}) := (D_{p_1},\dots,D_{p_n}) \in \mathbb{R}^n,
\]
and let
\[
  \delta_{p_i} = 1 - D_{p_i}, \qquad
  g_i := \abs{\delta_{p_i} - \delta_\ast}, \quad \delta_\ast := \frac{1}{137}.
\]
We endow \(\mathbb{R}^n\) with the usual \(\ell^\infty\) and \(\ell^2\) norms:
\[
  \|\mathbf{D}\|_\infty := \max_i |D_{p_i}|, \qquad
  \|\mathbf{D}\|_2 := \left(\sum_i D_{p_i}^2\right)^{1/2}.
\]

\subsection{Embedding depth as a bounded linear functional}

Define the embedding depth functional on \(\mathbb{R}^n\) by
\[
  \mathcal{E}(\mathbf{D}) := \frac{1}{n}\sum_{i=1}^n D_{p_i}.
\]

\begin{lemma}
As a linear functional \(\mathcal{E}\colon (\mathbb{R}^n,\|\cdot\|_\infty)\to\mathbb{R}\), we have
\[
  \|\mathcal{E}\|_{\mathrm{op}} = 1.
\]
As a map \(\mathcal{E}\colon (\mathbb{R}^n,\|\cdot\|_2)\to\mathbb{R}\), we have
\[
  \|\mathcal{E}\|_{\mathrm{op}} = \frac{1}{\sqrt{n}}.
\]
\end{lemma}

\begin{proof}
For the \(\ell^\infty\) case,
\[
  |\mathcal{E}(\mathbf{D})|
    = \frac{1}{n}\left|\sum_i D_{p_i}\right|
    \le \frac{1}{n} \sum_i |D_{p_i}|
    \le \frac{1}{n} \cdot n \|\mathbf{D}\|_\infty
    = \|\mathbf{D}\|_\infty.
\]
Taking the supremum over \(\mathbf{D}\neq 0\) shows \(\|\mathcal{E}\|_{\mathrm{op}}\le 1\), and equality is attained, e.g.\ by \(\mathbf{D} = (1,\dots,1)\).

For the \(\ell^2\) case, interpret \(\mathcal{E}(\mathbf{D}) = \langle \mathbf{D}, v\rangle\) with \(v=(1/n,\dots,1/n)\). Then
\[
  \|\mathcal{E}\|_{\mathrm{op}} = \|v\|_2 = \left(n\cdot \frac{1}{n^2}\right)^{1/2} = \frac{1}{\sqrt{n}}.
\]
\end{proof}

Thus embedding depth is a uniformly bounded linear observable on the configuration space, with norm at most one under reasonable choices of norm.

\subsection{Collapse index: lower and upper bounds}

Recall the definition
\[
  \mathrm{collapse}(\mathcal{C}) := \frac{1}{\min_{i<j}|D_{p_i} - D_{p_j}|},
\]
with the convention that for \(n=1\) the expression is undefined or set to \(+\infty\). Let
\[
  \Delta(\mathcal{C}) := \min_{i<j}|D_{p_i} - D_{p_j}|, \qquad n:=|\mathcal{C}|\ge 2.
\]

\begin{lemma}[Bounds in terms of gaps]
For any configuration \(\mathcal{C}\) with \(n\ge 2\),
\[
  0 < \Delta(\mathcal{C}) \le \|\mathbf{D}(\mathcal{C})\|_\infty,
\]
and
\[
  \mathrm{collapse}(\mathcal{C})
    = \frac{1}{\Delta(\mathcal{C})}
    \ge \frac{1}{\|\mathbf{D}(\mathcal{C})\|_\infty}.
\]
\end{lemma}

\begin{proof}
Since the \(D_{p_i}\) are distinct for distinct primes \(p_i\) under our construction, the minimum difference \(\Delta(\mathcal{C})\) is strictly positive. For any pair \(i\neq j\),
\[
  |D_{p_i} - D_{p_j}| \le |D_{p_i}| + |D_{p_j}| \le 2\|\mathbf{D}\|_\infty,
\]
so, in particular, \(\Delta(\mathcal{C}) \le 2\|\mathbf{D}\|_\infty\). The simpler bound \(\Delta(\mathcal{C}) \le \|\mathbf{D}\|_\infty\) holds because at least one \(D_{p_i}\) attains the maximum and differences cannot all exceed that maximum in magnitude.

Inverting the inequality \(\Delta(\mathcal{C}) \le \|\mathbf{D}\|_\infty\) yields
\[
  \mathrm{collapse}(\mathcal{C})
    = \frac{1}{\Delta(\mathcal{C})} \ge \frac{1}{\|\mathbf{D}\|_\infty}.
\]
\end{proof}

Thus, high collapse index necessarily implies large operator norm (in the sense of large \(\|\mathbf{D}\|_\infty\)) or extremely tight clustering of dimensions.

\subsection{Resonance score as a Lipschitz functional}

Define the resonance score
\[
  \mathrm{resonance}(\mathcal{C}) 
    := \frac{1}{n}\sum_{i=1}^n \frac{1}{1 + \lambda g_i},
\]
with \(\lambda = 10^6\) in the simulator.%
\footnote{This choice of scaling is purely numerical and may be treated as a parameter in analytical estimates. [code\_file:43]} 
Consider the map \(\Phi\colon \mathbb{R}^n \to \mathbb{R}^n\),
\[
  \Phi(\mathbf{D})_i = \frac{1}{1 + \lambda|\delta_{p_i} - \delta_\ast|},
\]
and the averaging operator \(\mathcal{A}(\mathbf{x}) = \frac{1}{n}\sum_i x_i\). Then
\[
  \mathrm{resonance}(\mathcal{C}) = (\mathcal{A}\circ\Phi)(\mathbf{D}(\mathcal{C})).
\]

\begin{lemma}[Lipschitz bound for resonance]
The map \(\mathbf{D}\mapsto \mathrm{resonance}(\mathbf{D})\) is globally Lipschitz on \(\mathbb{R}^n\) with respect to \(\ell^\infty\), with
\[
  \left|\mathrm{resonance}(\mathbf{D}) - \mathrm{resonance}(\mathbf{D}')\right|
  \le L \,\|\mathbf{D} - \mathbf{D}'\|_\infty,
\]
where 
\[
  L \le \frac{\lambda}{(1+\lambda m)^2},
\]
and \(m\) is a lower bound on all gaps \(g_i\) encountered along the line segment between \(\mathbf{D}\) and \(\mathbf{D}'\). In particular, for small gaps one can take the crude bound \(L \le \lambda\).
\end{lemma}

\begin{proof}
Fix \(i\). Consider the scalar function
\[
  f(d) := \frac{1}{1 + \lambda|\,(1-d) - \delta_\ast\,|},
\]
where \(d\) stands in for \(D_{p_i}\). Its derivative exists away from the kink at \(d = 1 - \delta_\ast\) and satisfies
\[
  |f'(d)| = \frac{\lambda}{(1 + \lambda|\, (1-d) - \delta_\ast\,|)^2}
  \le \frac{\lambda}{(1+\lambda m)^2}
\]
whenever \(|(1-d) - \delta_\ast|\ge m\). Thus \(f\) is Lipschitz with constant at most \(\lambda/(1+\lambda m)^2\) on any region where the gap is bounded below by \(m\). Since \(\mathcal{A}\) has operator norm \(1\) on \(\ell^\infty\), the composition \(\mathcal{A}\circ\Phi\) has Lipschitz constant bounded by the same quantity:
\[
  \left|\mathrm{resonance}(\mathbf{D}) - \mathrm{resonance}(\mathbf{D}')\right|
  \le \frac{\lambda}{(1+\lambda m)^2} \|\mathbf{D} - \mathbf{D}'\|_\infty.
\]
For the crude global bound, note that \(f\) is uniformly Lipschitz with constant at most \(\lambda\) (attained as the gap tends to zero), giving \(L\le \lambda\).
\end{proof}

This shows that small perturbations of the configuration dimensions translate into controlled perturbations of the resonance score, with an explicit Lipschitz constant depending on the scaling parameter and minimum gap.

\subsection{Mean gap as a semi–norm}

The mean absolute gap
\[
  \overline{g}(\mathcal{C}) := \frac{1}{n}\sum_{i=1}^n g_i
\]
is, up to the change of variables \(D_{p_i}\mapsto \delta_{p_i}\), a standard \(\ell^1\) semi–norm around a fixed center \(\delta_\ast\). It satisfies
\[
  0 \le \overline{g}(\mathcal{C}) \le \max_i g_i \le \|\bm{\delta}-\delta_\ast\mathbf{1}\|_\infty,
\]
and is a Lipschitz function of \(\mathbf{D}\) with Lipschitz constant \(1\) in \(\ell^\infty\).

\begin{lemma}
The map \(\mathbf{D}\mapsto \overline{g}(\mathbf{D})\) is 1–Lipschitz with respect to \(\ell^\infty\).
\end{lemma}

\begin{proof}
Consider again \(h(d) = |(1-d) - \delta_\ast|\). Clearly \(h\) is 1–Lipschitz on \(\mathbb{R}\). The mean gap is a composition of \(h\) applied coordinatewise and the averaging operator:
\[
  \overline{g}(\mathbf{D}) = \frac{1}{n}\sum_i h(D_{p_i}).
\]
A coordinatewise 1–Lipschitz map composed with an averaging operator of norm 1 in \(\ell^\infty\) yields a 1–Lipschitz map, so the claim follows.
\end{proof}

\nocite{*}
\bibliographystyle{unsrt}  
\bibliography{references}  


\end{document}
